% Options for packages loaded elsewhere
\PassOptionsToPackage{unicode}{hyperref}
\PassOptionsToPackage{hyphens}{url}
\PassOptionsToPackage{dvipsnames,svgnames,x11names}{xcolor}
\documentclass[
  11pt,
]{article}
\usepackage{xcolor}
\usepackage[margin=1in]{geometry}
\usepackage{amsmath,amssymb}
\setcounter{secnumdepth}{-\maxdimen} % remove section numbering
\usepackage{iftex}
\ifPDFTeX
  \usepackage[T1]{fontenc}
  \usepackage[utf8]{inputenc}
  \usepackage{textcomp} % provide euro and other symbols
\else % if luatex or xetex
  \usepackage{unicode-math} % this also loads fontspec
  \defaultfontfeatures{Scale=MatchLowercase}
  \defaultfontfeatures[\rmfamily]{Ligatures=TeX,Scale=1}
\fi
\usepackage{lmodern}
\ifPDFTeX\else
  % xetex/luatex font selection
\fi
% Use upquote if available, for straight quotes in verbatim environments
\IfFileExists{upquote.sty}{\usepackage{upquote}}{}
\IfFileExists{microtype.sty}{% use microtype if available
  \usepackage[]{microtype}
  \UseMicrotypeSet[protrusion]{basicmath} % disable protrusion for tt fonts
}{}
\makeatletter
\@ifundefined{KOMAClassName}{% if non-KOMA class
  \IfFileExists{parskip.sty}{%
    \usepackage{parskip}
  }{% else
    \setlength{\parindent}{0pt}
    \setlength{\parskip}{6pt plus 2pt minus 1pt}}
}{% if KOMA class
  \KOMAoptions{parskip=half}}
\makeatother
\usepackage{color}
\usepackage{fancyvrb}
\newcommand{\VerbBar}{|}
\newcommand{\VERB}{\Verb[commandchars=\\\{\}]}
\DefineVerbatimEnvironment{Highlighting}{Verbatim}{commandchars=\\\{\}}
% Add ',fontsize=\small' for more characters per line
\newenvironment{Shaded}{}{}
\newcommand{\AlertTok}[1]{\textcolor[rgb]{1.00,0.00,0.00}{\textbf{#1}}}
\newcommand{\AnnotationTok}[1]{\textcolor[rgb]{0.38,0.63,0.69}{\textbf{\textit{#1}}}}
\newcommand{\AttributeTok}[1]{\textcolor[rgb]{0.49,0.56,0.16}{#1}}
\newcommand{\BaseNTok}[1]{\textcolor[rgb]{0.25,0.63,0.44}{#1}}
\newcommand{\BuiltInTok}[1]{\textcolor[rgb]{0.00,0.50,0.00}{#1}}
\newcommand{\CharTok}[1]{\textcolor[rgb]{0.25,0.44,0.63}{#1}}
\newcommand{\CommentTok}[1]{\textcolor[rgb]{0.38,0.63,0.69}{\textit{#1}}}
\newcommand{\CommentVarTok}[1]{\textcolor[rgb]{0.38,0.63,0.69}{\textbf{\textit{#1}}}}
\newcommand{\ConstantTok}[1]{\textcolor[rgb]{0.53,0.00,0.00}{#1}}
\newcommand{\ControlFlowTok}[1]{\textcolor[rgb]{0.00,0.44,0.13}{\textbf{#1}}}
\newcommand{\DataTypeTok}[1]{\textcolor[rgb]{0.56,0.13,0.00}{#1}}
\newcommand{\DecValTok}[1]{\textcolor[rgb]{0.25,0.63,0.44}{#1}}
\newcommand{\DocumentationTok}[1]{\textcolor[rgb]{0.73,0.13,0.13}{\textit{#1}}}
\newcommand{\ErrorTok}[1]{\textcolor[rgb]{1.00,0.00,0.00}{\textbf{#1}}}
\newcommand{\ExtensionTok}[1]{#1}
\newcommand{\FloatTok}[1]{\textcolor[rgb]{0.25,0.63,0.44}{#1}}
\newcommand{\FunctionTok}[1]{\textcolor[rgb]{0.02,0.16,0.49}{#1}}
\newcommand{\ImportTok}[1]{\textcolor[rgb]{0.00,0.50,0.00}{\textbf{#1}}}
\newcommand{\InformationTok}[1]{\textcolor[rgb]{0.38,0.63,0.69}{\textbf{\textit{#1}}}}
\newcommand{\KeywordTok}[1]{\textcolor[rgb]{0.00,0.44,0.13}{\textbf{#1}}}
\newcommand{\NormalTok}[1]{#1}
\newcommand{\OperatorTok}[1]{\textcolor[rgb]{0.40,0.40,0.40}{#1}}
\newcommand{\OtherTok}[1]{\textcolor[rgb]{0.00,0.44,0.13}{#1}}
\newcommand{\PreprocessorTok}[1]{\textcolor[rgb]{0.74,0.48,0.00}{#1}}
\newcommand{\RegionMarkerTok}[1]{#1}
\newcommand{\SpecialCharTok}[1]{\textcolor[rgb]{0.25,0.44,0.63}{#1}}
\newcommand{\SpecialStringTok}[1]{\textcolor[rgb]{0.73,0.40,0.53}{#1}}
\newcommand{\StringTok}[1]{\textcolor[rgb]{0.25,0.44,0.63}{#1}}
\newcommand{\VariableTok}[1]{\textcolor[rgb]{0.10,0.09,0.49}{#1}}
\newcommand{\VerbatimStringTok}[1]{\textcolor[rgb]{0.25,0.44,0.63}{#1}}
\newcommand{\WarningTok}[1]{\textcolor[rgb]{0.38,0.63,0.69}{\textbf{\textit{#1}}}}
\usepackage{longtable,booktabs,array}
\usepackage{calc} % for calculating minipage widths
% Correct order of tables after \paragraph or \subparagraph
\usepackage{etoolbox}
\makeatletter
\patchcmd\longtable{\par}{\if@noskipsec\mbox{}\fi\par}{}{}
\makeatother
% Allow footnotes in longtable head/foot
\IfFileExists{footnotehyper.sty}{\usepackage{footnotehyper}}{\usepackage{footnote}}
\makesavenoteenv{longtable}
\setlength{\emergencystretch}{3em} % prevent overfull lines
\providecommand{\tightlist}{%
  \setlength{\itemsep}{0pt}\setlength{\parskip}{0pt}}
\usepackage{bookmark}
\IfFileExists{xurl.sty}{\usepackage{xurl}}{} % add URL line breaks if available
\urlstyle{same}
\hypersetup{
  colorlinks=true,
  linkcolor={Maroon},
  filecolor={Maroon},
  citecolor={Blue},
  urlcolor={Blue},
  pdfcreator={LaTeX via pandoc}}

\author{}
\date{}

\begin{document}

\section{Finite Distinction Carriers and History
Escape}\label{finite-distinction-carriers-and-history-escape}

\textbf{Subtitle:} Observational equivalence and minimal deterministic
quotients

\textbf{Author:} Jeremy H. Carroll

\textbf{Date:} 2026-10-08

\textbf{Version:} 1.1.0

\textbf{Status:} Author-review prose preprint; not an external
publication

\textbf{License:} Creative Commons Attribution 4.0 International
(CC-BY-4.0)

\textbf{AuthorityEffect:} None

\subsection{Abstract}\label{abstract}

For a finite deterministic system with a fixed observation map, two
states may have the same present observation but different observations
after a sequence of updates. We characterize this history-dependent
distinction by a descending sequence of equivalence relations. The
recursive refinement equals comparison over all input words up to a
given length. It stabilizes within the number of states minus the number
of initial observation classes; every pair removed at a refinement step
has a shortest separating word of the corresponding length. The
stabilized relation is the largest observation-preserving update
congruence, and its quotient is the smallest deterministic quotient
preserving every state's future observations. We also give a sufficient
report-factorization condition, its proof, and a counterexample when
reports introduce an additional observation. The arguments, including
empty-state and empty-alphabet cases, are written in the text. In this
package only Theorem WD2-T1 is a Lean 4 theorem. Lemma WD2-L1 and
records WD2-T2 through WD2-T6 are finite-instance checks or prose
arguments and are not labeled proved. The results are an exposition of
finite-state observational equivalence and classical state minimization,
not a new reconstruction of physical theory.

\subsection{1. Question and relation to state
minimization}\label{question-and-relation-to-state-minimization}

An initially hidden state difference can become visible after specified
updates. This differs from processing an unchanged observation: the
observer is allowed to update the state and then observe it again. The
question is which initially indistinguishable pairs can be separated in
this way, and how long a shortest separating sequence can be.

The model is a finite deterministic Moore machine with a finite output
alphabet. Its refinement and quotient results belong to classical
automata theory {[}1,2{]}. We claim no novelty for partition refinement,
observational equivalence, or minimal deterministic quotients. The
purpose is to state their hypotheses, witness order, and observation
boundary in one finite formulation. The arguments below are
self-contained; no unpublished companion result or unavailable
proof-assistant module is required.

\textbf{Package status (2026-10-09).} Theorem WD2-T1 is elaborated in
Lean 4. Reflexivity, symmetry, and transitivity are
\texttt{Wd2.related\_refl}, \texttt{Wd2.related\_symm}, and
\texttt{Wd2.related\_trans}. Inclusion and stability are
\texttt{Wd2.related\_mono} and \texttt{Wd2.related\_stable}. The
five-state chain at horizon 3 is \texttt{Wd2.sharp\_m5\_bound}. The
general class-count bound, the no-escape equivalence, the shortest-word
construction, the quotient mapping property, and the
report-factorization argument are not Lean theorems in this package. The
Python checker evaluates the sharp chain, the update-stable pair, and
the factored report.

\subsection{2. Finite model and word
convention}\label{finite-model-and-word-convention}

Fix finite sets \(S\) of states and \(\Sigma\) of input symbols, an
integer \(K\ge0\), maps \[
 u:S\times\Sigma\longrightarrow S,\qquad
 p:S\longrightarrow\{0,1\}^K,
\] and write \(u_\sigma(s)=u(s,\sigma)\). The binary output
representation is a choice, not an assertion that binary distinctions
are primitive. The proofs also apply to a map into any finite output
set.

Words are applied \textbf{left to right}. With \(\varepsilon\) denoting
the empty word, define \[
 u_\varepsilon=\operatorname{id}_S,\qquad
 u_{w\sigma}=u_\sigma\circ u_w.
\] Induction on the length of \(v\) gives \(u_{wv}=u_v\circ u_w\), and
in particular \(u_{\sigma v}=u_v\circ u_\sigma\). This convention
matters when updates do not commute. No order on the alphabet is
assumed.

For a map \(f:S\to Y\), its kernel relation is
\(\ker f=\{(s,t):f(s)=f(t)\}\). An \emph{update congruence} is an
equivalence relation \(E\) on \(S\) such that \(sEt\) implies
\(u_\sigma(s)E u_\sigma(t)\) for every symbol \(\sigma\). It preserves
the declared observation when \(E\subseteq\ker p\).

\subsection{3. Finite-horizon
observations}\label{finite-horizon-observations}

Set \(R_0=\ker p\) and, for integers \(n\ge0\), define \[
 R_{n+1}=R_0\cap\bigcap_{\sigma\in\Sigma}
              (u_\sigma\times u_\sigma)^{-1}(R_n).
\] The intersection over an empty alphabet is \(S\times S\). An
initially ambiguous pair \((s,t)\in R_0\) \emph{escapes at depth}
\(n\ge1\) if \((s,t)\in R_{n-1}\setminus R_n\). A \emph{separating word}
satisfies \(p(u_w(s))\ne p(u_w(t))\); an \emph{escape witness} is a
separating word of shortest length. ``Shortest'' does not mean a unique
word. A tie-breaking choice would require an additional order on
\(\Sigma\).

\textbf{Lemma WD2-L1 (word characterization).} For every integer
\(n\ge0\), \[
 R_n=\bigcap_{|w|\le n}\ker(p\circ u_w).
\]

\emph{Proof.} At \(n=0\) the only word is \(\varepsilon\), so the
formula is the definition of \(R_0\). Every nonempty word of length at
most \(n+1\) has a first symbol and a suffix: \(w=\sigma v\),
\(|v|\le n\). By the induction hypothesis, equality of all its readouts
is equivalent to \((u_\sigma(s),u_\sigma(t))\in R_n\) for every
\(\sigma\). The empty word additionally requires \((s,t)\in R_0\). These
are exactly the recursive conditions for \(R_{n+1}\). \(\square\)

\subsection{4. Refinement and
stabilization}\label{refinement-and-stabilization}

\textbf{Theorem WD2-T1 (refinement and fixed-point law).} Every \(R_n\)
is an equivalence relation and \(R_{n+1}\subseteq R_n\). If
\(R_{n+1}=R_n\), then \(R_m=R_n\) for every \(m\ge n\).

\emph{Proof.} Kernels of maps are equivalence relations, and
intersections preserve that property. Lemma L1 therefore establishes
equivalence. The set of words of length at most \(n\) is contained in
the set for \(n+1\), establishing the inclusion. Define \[
 F(E)=R_0\cap\bigcap_{\sigma\in\Sigma}
                     (u_\sigma\times u_\sigma)^{-1}(E).
\] If \(R_{n+1}=R_n\), then \(F(R_n)=R_n\). Applying \(F\) repeatedly
gives the claimed equality at all later indices. \(\square\)

\textbf{Classical record WD2-T2 (finite stabilization and a counting bound).} Not a theorem of this package. Put
\(b_0=|S/R_0|\). There is a least integer \(N\ge0\) such that
\(R_N=R_{N+1}\), and \[
 N\le |S|-b_0,\qquad R_\infty:=\bigcap_{n\ge0}R_n=R_N.
\] The number \(N\) counts strict refinement rounds. An implementation
can need one additional comparison to establish \(R_N=R_{N+1}\).

\emph{Proof.} A strict inclusion of finite equivalence relations splits
at least one class, increasing the number of classes by at least one.
Starting with \(b_0\) classes, no more than \(|S|-b_0\) consecutive
strict inclusions can occur. The first equality is permanent by T1. The
descending intersection is therefore its stabilized value. If \(S\) is
empty, both class counts and \(N\) are zero. If \(\Sigma\) is empty, the
recursion is constant and \(N=0\). \(\square\)

For sharpness when \(m=|S|\ge2\), take \(S=\{0,\ldots,m-1\}\), a single
symbol \(a\), \(u_a(i)=\min(i+1,m-1)\), and \(p(i)=1\) precisely for
\(i=m-1\). At horizon \(j\le m-2\), the lowest class is
\(\{0,\ldots,m-2-j\}\) and all larger states are singleton classes.
Indeed state \(i<m-1\) first reads one after exactly \(m-1-i\) updates.
Thus \(b_0=2\) and \(N=m-2=|S|-b_0\) (including \(N=0\) when \(m=2\)).

\subsection{5. No escape and shortest
witnesses}\label{no-escape-and-shortest-witnesses}

\textbf{Classical record WD2-T3 (no-escape criterion).} Not a theorem of this package. The following are
equivalent:

\begin{enumerate}
\def\labelenumi{\arabic{enumi}.}
\tightlist
\item
  \(R_0\) is an update congruence;
\item
  \(R_1=R_0\);
\item
  \(R_n=R_0\) for every \(n\ge0\);
\item
  no pair in \(R_0\) has a separating word.
\end{enumerate}

\emph{Proof.} Since \(R_1\subseteq R_0\), equality holds exactly when
every pair in \(R_0\) is carried by every update back into \(R_0\). This
is condition 1. Condition 2 implies 3 by T1, and 3 implies 2 by taking
\(n=1\). By L1, condition 3 says that every initially equivalent pair
has equal readouts after every finite word, which is exactly condition
4. \(\square\)

\textbf{Classical record WD2-T4 (shortest separating-word construction).} Not a theorem of this package. If
\((s,t)\in R_n\setminus R_{n+1}\), then its shortest separating-word
length is \(n+1\). Such a word can be constructed as a first symbol
followed by a length-\(n\) suffix. Consequently every escaping pair has
escape depth at most \(N\le |S|-b_0\).

\emph{Proof.} Membership in \(R_n\) excludes separation by any word of
length at most \(n\), by L1. Failure of \(R_{n+1}\), together with
membership in \(R_0\), provides a symbol \(\sigma\) for which the
successor pair lies outside \(R_n\). For \(n=0\), the empty suffix
separates that successor pair. For \(n\ge1\), the recursion for
membership of the original pair in \(R_n\) also places the successor
pair in \(R_{n-1}\). Induction supplies a shortest separating suffix
\(v\) of length \(n\). The word \(\sigma v\) separates the original pair
because \(u_{\sigma v}=u_v\circ u_\sigma\). Its length is \(n+1\), which
is minimal by the first observation. No pair leaves the sequence after
its stabilization at \(R_N\), giving the depth bound. \(\square\)

\subsection{6. The complete finite behavioral
quotient}\label{the-complete-finite-behavioral-quotient}

\textbf{Classical record WD2-T6 (largest congruence and smallest deterministic quotient).} Not a theorem of this package.
The relation \(R_\infty\) is the largest update congruence
contained in \(\ker p\). The quotient \(Q_*=S/R_\infty\) has
well-defined updates and output \[
 \bar u_\sigma([s])=[u_\sigma(s)],\qquad \bar p([s])=p(s).
\] Let \(q:S\twoheadrightarrow Q\) be any surjection with deterministic
maps \(v_\sigma:Q\to Q\) and \(o:Q\to\{0,1\}^K\) satisfying \[
 q\circ u_\sigma=v_\sigma\circ q,\qquad p=o\circ q.
\] There is a unique surjection \(h:Q\twoheadrightarrow Q_*\) such that
\(\pi=h\circ q\), where \(\pi(s)=[s]\). It commutes with updates and
preserves output. In particular \(|Q|\ge|Q_*|\). Equality makes \(h\) an
isomorphism of the deterministic output systems.

\emph{Proof.} Stabilization and the recursion show that \(R_\infty\) is
an update congruence contained in \(R_0\). If \(E\subseteq R_0\) is any
update congruence, induction on word length shows that \(sEt\) implies
\(u_w(s)E u_w(t)\) for every word. Their observations agree, and L1
gives \(E\subseteq R_n\) for every \(n\), hence \(E\subseteq R_\infty\).
This establishes maximality and also makes the displayed quotient maps
well-defined.

For a specified quotient \(q\), its kernel is an equivalence relation.
The update equation makes it an update congruence, and the output
equation places it inside \(R_0\). Maximality yields
\(\ker q\subseteq R_\infty\). Define \(h(q(s))=[s]\). Kernel inclusion
makes this independent of the choice of \(s\), and surjectivity of \(q\)
gives existence and uniqueness on all of \(Q\). Every class \([s]\) lies
in the image, so \(h\) is onto. Moreover, \[
 h(v_\sigma(q(s)))=[u_\sigma(s)]
      =\bar u_\sigma(h(q(s))),
\] and \(o(q(s))=p(s)=\bar p(h(q(s)))\). This establishes compatibility.
Finite surjectivity gives the cardinal bound; at equality \(h\) is
bijective, and its inverse also respects the maps. \(\square\)

Equivalently, the \emph{behavior} of a state is the function
\(\beta_s:\Sigma^*\to\{0,1\}^K\) given by \(\beta_s(w)=p(u_w(s))\).
Lemma L1 identifies \(R_\infty\) with equality of behaviors. Thus
\(Q_*\) classifies every state by all its future readouts, and the
induced update sends a behavior \(\beta_s\) to
\(w\mapsto\beta_s(\sigma w)\). This completes the classification of
states for the specified finite model. It does not classify physical
theories or select quantum mechanics.

The cardinality assertion concerns quotients preserving readouts from
\textbf{every} state in \(S\). A model restricted to the states
reachable from one initial state is a different minimization problem. If
\(S\) is empty, a surjective \(q\) has empty codomain, and the statement
still holds.

\subsection{7. Reports are additional observations unless
constrained}\label{reports-are-additional-observations-unless-constrained}

Fix a set \(\Lambda\) and a report \(r:S\times\Sigma\to\Lambda\). For a
word \(w=\sigma_1\cdots\sigma_m\), set \(s_0=s\) and
\(s_i=u_{\sigma_i}(s_{i-1})\). The report at step \(i\) is
\(r(s_{i-1},\sigma_i)\): it is attached to the update from \(s_{i-1}\)
to \(s_i\). This specifies the timing convention.

\textbf{Classical record WD2-T5 (sufficient report-factorization condition).}
Not a theorem of this package. Suppose there is a map
\(g:\{0,1\}^K\times\Sigma\times\{0,1\}^K\to\Lambda\) such that \[
 r(s,\sigma)=g(p(s),\sigma,p(u_\sigma(s))).
\] If \((s,t)\in R_n\), then their report sequences agree along every
common word of length at most \(n\). Consequently adjoining these report
sequences to all profile histories through horizon \(n\) leaves the
relation \(R_n\) unchanged.

\emph{Proof.} For each prefix of length \(i\le n\), L1 gives equality of
the profile values reached from \(s\) and \(t\). In particular the
profiles before and after step \(i\) agree, and the input symbol is the
same. Applying \(g\) therefore gives equal reports at that step. This
holds at every step of every such word. The profile comparisons already
impose \(R_n\); the equal reports impose no further condition.
\(\square\)

This is sufficient, not a necessary characterization. At horizon
infinity a report constant on \(R_\infty\)-classes for each input can
also preserve behavioral equivalence without factoring through just the
two endpoint profiles. For example, let states \(x,y,z\) have profiles
\(0,0,1\), let \(a\) act identically, and let \(b\) send \(x\) to \(x\)
and both \(y,z\) to \(z\). Here \(R_\infty\) is equality. A report
distinguishing \(x,y\) on input \(a\) preserves that equality relation
but cannot factor through the identical endpoint profile pair \((0,0)\)
for \(a\).

Without any report restriction even the one-step conclusion fails. Take
\(S=\{0,1\}\), one identity update \(a\), constant profile \(p=0\), and
\(r(s,a)=s\). Every profile history agrees, but the first reports
differ. The report has supplied an additional observation, not recovered
a difference from the unchanged profile.

\subsection{8. Worked witnesses and boundary
cases}\label{worked-witnesses-and-boundary-cases}

\textbf{Three-state escape.} Let \(S=\{x,y,z\}\), \(\Sigma=\{a\}\),
\(p(x)=p(y)=0\), \(p(z)=1\), and \(u_a(x)=x\), \(u_a(y)=u_a(z)=z\).
Initially the classes are \(\{x,y\}\) and \(\{z\}\). One refinement
makes every state a singleton; \(a\) is a shortest witness for \(x,y\).
Thus \(N=1=3-2\).

\textbf{Noncommuting updates.} Let \(S=\{0,1,2,3\}\), profiles
\((0,0,0,1)\), and specify the update values in state order by \[
 u_a=(1,2,2,3),\qquad u_b=(0,0,3,3).
\] States \(0,1\) have equal readouts at the empty word and after either
one-letter word. Under \(ab\) they reach \(0,3\), whose profiles differ.
Under \(ba\) they both reach \(1\). Thus \(ab\) is a shortest separating
word of length two, and reversing the word destroys this witness. In
fact \(R_1\) has classes \(\{0,1\},\{2\},\{3\}\) and \(R_2\) is
equality.

\textbf{No symbols or no output bits.} If \(\Sigma\) is empty, the only
word is \(\varepsilon\) and no initially ambiguous pair escapes. If
\(K=0\), the output set is a singleton, so all pairs have equal profiles
forever, regardless of the updates. For empty \(S\), every relation and
quotient is empty and every assertion about a state pair is vacuous.

\subsection{Mechanism inventory}\label{mechanism-inventory}

\begin{longtable}[]{@{}
  >{\raggedright\arraybackslash}p{(\linewidth - 2\tabcolsep) * \real{0.5000}}
  >{\raggedright\arraybackslash}p{(\linewidth - 2\tabcolsep) * \real{0.5000}}@{}}
\toprule\noalign{}
\begin{minipage}[b]{\linewidth}\raggedright
Object
\end{minipage} & \begin{minipage}[b]{\linewidth}\raggedright
Mathematical role
\end{minipage} \\
\midrule\noalign{}
\endhead
\bottomrule\noalign{}
\endlastfoot
Finite \(S\) & States being compared \\
Finite \(\Sigma\) and \(u_\sigma:S\to S\) & Permitted deterministic
updates \\
Finite profile \(p:S\to\{0,1\}^K\) & Declared observation \\
Equivalence relations \(R_n\) & Equality of readouts through horizon
\(n\) \\
Separating word & Finite witness of observable difference \\
\(S/R_\infty\) & Smallest deterministic quotient preserving all states'
future readouts \\
Optional report \(r\) & Additional observation with explicit timing and
factorization hypothesis \\
\textbf{AuthorityEffect} & None \\
\end{longtable}

\subsection{Verification scope}\label{verification-scope}

Theorem WD2-T1 is elaborated in Lean 4 with no \texttt{sorry}. Lemma
WD2-L1 and records WD2-T2 through WD2-T6 have prose arguments above.
The Python checker evaluates finite instances of those statements; their
general forms are not Lean theorems in this package. From the package
root, this command checks pinned files, those finite tests, and the Lean
build:

\begin{Shaded}
\begin{Highlighting}[]
\ExtensionTok{py} \AttributeTok{{-}3.12}\NormalTok{ verify\_deposit.py }\AttributeTok{{-}{-}lean} \AttributeTok{{-}{-}test}\NormalTok{ code/tests/test\_proved\_rows.py}
\end{Highlighting}
\end{Shaded}

The reading notebook contains prose only and requires no kernel
execution.

No cochain complex, Hilbert space, repair norm or physical probability
rule is defined here. A quotient by behavioral equivalence is a quotient
of a finite set, not a Banach quotient. Nothing here derives memory,
agency, consciousness, or quantum theory from a deterministic update
system.

\subsection{References}\label{references}

\begin{enumerate}
\def\labelenumi{\arabic{enumi}.}
\tightlist
\item
  L. Trevisan, \emph{CS103: Mathematical Foundations of Computing},
  Lecture 16, state minimization, Stanford University (2010).
  \url{https://theory.stanford.edu/~trevisan/cs103-10/lecture16.pdf}.
\item
  J. E. Hopcroft, R. Motwani, and J. D. Ullman, \emph{Introduction to
  Automata Theory, Languages, and Computation}, third edition, Pearson
  (2007), Chapter 4. ISBN 978-0-321-45536-9.
\end{enumerate}

\end{document}
